\clearpage\fancyhead[L]{\small GAUSS AN BODE.}

\begin{center}\textbf{GAUSS AN BODE.}\par\vspace{4pt}\hrule\vspace{4pt}\fontsize{10}{12}\selectfont Astronomisches Jahrbuch für 1813. Seite 236\,..\,238. Berlin 1810.\par\vspace{4pt}\hrule\vspace{4pt}\end{center}
\begin{flushright}Göttingen, den 2. September 1810.\end{flushright}

\par\indent ---\enspace---\enspace---

\par\indent Die Elemente, nach denen diese Ephemeride [der Pallas für 1810 Sept. 16 bis 1811 Juli 1 von \textsc{Gauss}] berechnet wurde, beruhen nur auf den letzten 4 Oppositionen und werden ausreichen, um den Planeten diesmal aufzufinden. Ich habe aber auch schon vor einiger Zeit die Elemente berechnet, die möglichst genau mit allen 6 bisher beobachteten Oppositionen übereinstimmen und sich jetzt am besten als Grundlage für die Berechnung der Störungen eignen. Die Störungen beginnen inzwischen, sehr deutlich hervorzutreten. Ich werde unserer Gesellschaft der Wissenschaften demnächst eine Abhandlung über diese Rechnungen vorlegen. Sie werden bereits wissen, dass unsere Regierung sehr viel für die Astronomie tun will 200000 Franken sollen für den Bau der neuen Sternwarte verwendet werden. Professor \textsc{Harding} hat für eine Reise nach Paris eine Zuwendung von 4000 Franken erhalten, um dort die Lücken der Histoire céleste zu füllen und seine Sternkarten zu vervollständigen.

Die Ehre meiner kürzlichen Aufnahme in Ihre Akademie der Wissenschaften weiß ich sehr zu schätzen.
